\section{第24章：双重差分的发展与问题}
\subsection{面板数据与双向固定效应}
\sourcepages{Chapter24-25.pdf，第2、1页；第1页接续第2页底部的假设。}
\textbf{面板数据：}在不同时点重复观测同一个单位。例子同时比较不同时间、不同地区的单位。图中有早期受处理组（6月1日）、晚期受处理组（7月15日）和对照组；另有一条较浅的轨迹，标为\emph{TWFE预测／反事实}。
\begin{center}
\begin{tikzpicture}[x=1.1cm,y=0.68cm,>=stealth]
\draw[->] (0,0)--(9,0) node[right]{时间};
\draw[->] (0,0)--(0,3.8) node[above]{结果};
\draw[dashed] (3,0)--(3,3.8);
\draw[dashed] (5.6,0)--(5.6,3.8);
\node[below] at (3,0){6月1日}; \node[below] at (5.6,0){7月15日};
\draw (0.5,.4)--(2.9,.7)--(5.5,.9)--(7.6,1.1) node[right]{对照组};
\draw (0.5,1)--(3,1.3)--(5.6,1.5)--(5.8,1.85)--(7.6,2) node[right]{7月15日};
\draw (0.5,2)--(3,2.25)--(3.2,3)--(7.6,3.3) node[right]{6月1日};
\draw[gray,dashed] (3,2.25)--(7.6,2.6);
\end{tikzpicture}

{\small 按手绘轨迹重建的示意图，不代表具体数值数据。}
\end{center}
按干预开始日期 \(g\) 给单位分组（例中分为三组）。
\notearrow{black}{干预开始日期}{队列标签 \(g\)} 笔记写出的目标为
\[
\operatorname{ATT}=E[Y_1\mid t>g]-E[Y_0\mid t>g].
\]
\begin{reviewnote}{qualification}{B2-003}{明确队列—时间目标量}{Chapter24-25.pdf, pp. 2}
更精确的定义是 \(\operatorname{ATT}(g,t)=E[Y_t(g)-Y_t(\infty)\mid G=g]\)，其中 \(g\) 是受处理队列，\(t\) 是处理后的时点。还需说明处理从 \(t=g\) 开始，还是只在其后开始。总体ATT还需要明确聚合权重。
\par\smallskip{\small\textit{参考：} \href{https://matheusfacure.github.io/python-causality-handbook/24-The-Diff-in-Diff-Saga.html}{原始资料1}.}
\end{reviewnote}

模型为
\[
\mathrm{Install}_{it}=\tau D_{it}+\gamma_i+\theta_t+e_{it}.
\]
这里 \(i\) 表示单位，\(t\) 表示时间，\(e_{it}\) 为噪声；\(\gamma_i\) 对应单位，\(\theta_t\) 对应时间。\inkblue{用OLS估计，得到\textbf{双向固定效应（TWFE）}模型。}

\paragraph{我的批注。}
\inkblue{同时控制所有特征。}
\notearrow{red}{两组、两期}{TWFE就是DiD。}\begin{reviewnote}{original}{B2-002}{固定效应究竟控制了什么}{Chapter24-25.pdf，第2页}
单位固定效应吸收单位之间不随时间变化的差异；时间固定效应吸收某一时点对所有单位共同发生的冲击。它们不能控制任意的、因单位而异且随时间变化的混杂因素。可加性与识别假设仍然重要。
\par\smallskip{\small\textit{参考：} \href{https://matheusfacure.github.io/python-causality-handbook/24-The-Diff-in-Diff-Saga.html}{原始资料1}.}
\end{reviewnote}

\notearrow{black}{TWFE结果模型}{用同一组固定效应对结果与处理分别取残差。}
笔记将残差化后的回归写为
\[
\widetilde{\mathrm{Install}}_{it}=\tau\widetilde D_{it}+e_{it},
\]
对应的去均值表达式为：
\begin{align*}
\widetilde{\mathrm{Install}}_{it}
 &=\mathrm{Install}_{it}
 -\underbrace{\frac1T\sum_{t=0}^{T}\mathrm{Install}_{it}}_{\textcolor{red}{\text{时间平均}}}
 -\underbrace{\frac1N\sum_{i=0}^{N}\mathrm{Install}_{it}}_{\textcolor{red}{\text{单位平均}}},\\
\widetilde D_{it}
 &=D_{it}-\frac1T\sum_{t=0}^{T}D_{it}-\frac1N\sum_{i=0}^{N}D_{it}.
\end{align*}
\begin{reviewnote}{original}{B2-004}{双重去均值需要加回总体均值}{Chapter24-25.pdf, pp. 2, 6}
对平衡面板，令 \(i=1,\ldots,N\)、\(t=1,\ldots,T\)，应使用
\[
\widetilde z_{it}=z_{it}-\bar z_{i\cdot}-\bar z_{\cdot t}+\bar z_{\cdot\cdot},\qquad z\in\{Y,D\}.
\]
其中 \(\bar z_{i\cdot}=T^{-1}\sum_{s=1}^Tz_{is}\)、\(\bar z_{\cdot t}=N^{-1}\sum_{j=1}^Nz_{jt}\)、\(\bar z_{\cdot\cdot}=(NT)^{-1}\sum_{j,s}z_{js}\)。加回总体均值后，各行、各列均值才都为零。遗漏它只会使每个变换后的变量平移一个常数；若最后的回归另含截距，斜率不变。但这仍不是精确的固定效应残差，且无截距回归不一定给出相同斜率。非平衡面板应使用实际的固定效应投影，而不能直接套用这个简单公式。
\par\smallskip{\small\textit{参考：} \href{https://matheusfacure.github.io/python-causality-handbook/24-The-Diff-in-Diff-Saga.html}{原始资料1}, \href{https://matheusfacure.github.io/python-causality-handbook/25-Synthetic-Diff-in-Diff.html}{原始资料2}.}
\end{reviewnote}


\subsection{失效：处理效应异质性带来的问题}
效应异质时，TWFE经常有偏。没有时间效应的固定效应模型写为
\[
y_{it}=\tau D_{it}+\gamma_i+e_{it}.
\]
笔记列出的假设包括：
\begin{enumerate}
\item \textbf{函数形式：}效应不随时间异质变化（常数效应）；协变量线性进入；固定效应可加。
\item \textbf{严格外生性：}平行趋势；无预期效应；无未观测的时变混杂；过去的处理不影响当前结果（无延续效应）；过去的结果不影响当前处理（无反馈）。
\end{enumerate}
第二项接在左页，由主页面底部的箭头连接。
\notearrow{black}{第2页底部的严格外生性列表}{第1页接续：预期、混杂、延续效应和反馈。}\begin{reviewnote}{qualification}{B2-005}{严格外生性与DiD假设}{Chapter24-25.pdf，第1、2页}
静态固定效应回归中的一个标准条件是 \(E[e_{it}\mid D_{i1},\ldots,D_{iT},\gamma_i]=0\)（同时也以其他回归变量为条件）。平行趋势和无预期效应是DiD使用的潜在结果限制，并不与整张列表互相等价。通过改变模型与目标量，可以允许动态处理效应。
\par\smallskip{\small\textit{参考：} \href{https://matheusfacure.github.io/python-causality-handbook/24-The-Diff-in-Diff-Saga.html}{原始资料1}.}
\end{reviewnote}


\sourcepages{Chapter24-25.pdf, p. 3.}
接上：假设治疗效果存在时间异质性，而且反应是渐进的，例如从 \(0.1,0.2,\ldots\) 开始，每天增加。这时TWFE估计的ATT会有偏。
图中早处理组记为 \(1\)，晚处理组记为 \(2\)，从未处理组记为 \(3\)。把TWFE分解成多个 \(2\times2\) DiD，可得：
\[
1-3:\ \checkmark,\quad
2-3:\ \checkmark,\quad
1-2:\ \checkmark,\quad
2-1:\ \text{??}
\]
每个比较的第二项是对照。\textbf{我的直觉：}当早期组 \(1\) 被用作对照时，模型以为未处理结果 \(Y_0\) 会沿该组先前的趋势平行上升。\notearrow{blue}{估计出的未处理轨迹}{蓝色曲线是晚处理组的 \(\widehat Y_0\)。}
\begin{center}
\begin{tikzpicture}[x=.85cm,y=.65cm,>=stealth]
\draw[->] (0,0)--(8,0) node[right]{时间}; \draw[->] (0,0)--(0,3.4);
\draw[dashed] (2,0)--(2,3.2); \draw[dashed] (5,0)--(5,3.2);
\draw (0.5,1.2)--(2,1.4)--(3,2.6)--(7,3) node[right,yshift=4pt]{1：早处理};
\draw (0.5,.6)--(5,1.1)--(6,2.1)--(7,2.7) node[right,yshift=-4pt]{2：晚处理};
\draw[blue] (.5,.35)--(5,.85)--(7,2.2);
\node[blue,below] (b) at (6,-.5){估计的 \(Y_0\)};
\draw[->,blue] (b.north)--(6,1.55);
\node[left] at (0,2.8){\(2-1\)?};
\end{tikzpicture}
\end{center}
此示意图恢复“比较方式指向曲线”的关系；曲线高度仅用于说明。

笔记将下面的简写归于Goodman--Bacon（2019）：
\[
\operatorname*{plim}_{n\to\infty}\widehat\tau^{\mathrm{TWFE}}
 =\mathrm{VWATT}-\Delta\mathrm{ATT}.
\]
\notearrow{blue}{Bias：ATT随时间变化的程度}{\(\Delta\mathrm{ATT}\)}
\inkblue{若效果随时间增加（如图），Bias向下；反之亦然。\textbf{我的提醒：}当效果为负且随时间递减，即变得更负时，估计甚至会出现正负号错误！}\begin{reviewnote}{qualification}{B2-006}{偏差方向的说法针对图示比较}{Chapter24-25.pdf，第3页}
图解释了把已经受处理的单位当作对照造成的污染。总体TWFE偏差的方向和大小取决于各队列的效应、动态变化以及回归权重；该示意不是普遍成立的偏差符号定理。
\par\smallskip{\small\textit{参考：} \href{https://matheusfacure.github.io/python-causality-handbook/24-The-Diff-in-Diff-Saga.html}{原始资料1}.}
\end{reviewnote}


\subsection{事件研究设计}
用事件时间指标改写TWFE。笔记中的提前项／滞后项展开如下（保留其系数与下标标签）：
\begin{align*}
Y_{it}={}&\tau^{-K}D_{ik}^{<-K}
 +\sum_{k=-K}^{-2}\tau^{\mathrm{lead}}D_{it}^{k}
 +\sum_{k=0}^{L}\gamma_k^{\mathrm{lag}}D_{it}^{k}\\
 &+\gamma_k^{L+}D_{it}^{>L}+\gamma_i+e_{it},\\
D_{it}^{k}={}&\mathbf1\{t-\mathrm{Cohort}_i=k\}.
\end{align*}
\notearrow{black}{\(D_{it}^{k}=\mathbf1\{t-\mathrm{Cohort}_i=k\}\)}{相对时期为 \(k\) 时取1，否则取0。}
\notearrow{blue}{加入衡量单位距离开始受处理时期多远的变量}{事件时间指标 \(D_{it}^{k}\)。}
这里标为\emph{部分改进}。\begin{reviewnote}{original}{B2-007}{完整的提前项／滞后项回归}{Chapter24-25.pdf，第3页}
省略基准期 \(k=-1\) 的常见设定为
\begin{align*}
Y_{it}={}&\alpha_i+\theta_t+\beta_{<-K}\mathbf1\{t-G_i<-K\}\\
&+\sum_{k=-K}^{-2}\beta_k\mathbf1\{t-G_i=k\}
+\sum_{k=0}^{L}\beta_k\mathbf1\{t-G_i=k\}\\
&+\beta_{>L}\mathbf1\{t-G_i>L\}+e_{it}.
\end{align*}
从未受处理单位的所有事件时间指标均为零。原稿在单位效应前的分隔符不清楚，此处读作加号；不一致的下标则属于原稿记号问题。
\par\smallskip{\small\textit{参考：} \href{https://matheusfacure.github.io/python-causality-handbook/24-The-Diff-in-Diff-Saga.html}{原始资料1}.}
\end{reviewnote}
\begin{reviewnote}{qualification}{B2-008}{仅加入事件时间指标仍不能解决效应异质性}{Chapter24-25.pdf, pp. 3}
在交错采用且效应异质时，一个提前项或滞后项系数可能混入其他相对时期的效应，表面上的处理前趋势也可能来自效应异质性。若要作因果解释，应使用明确构造有效队列比较的估计方法。
\par\smallskip{\small\textit{参考：} \href{https://arxiv.org/abs/1804.05785}{原始资料1}.}
\end{reviewnote}


\subsection{改进：更灵活的函数形式}
\sourcepages{Chapter24-25.pdf, p. 4.}
新模型给每个个体在每个时点分配一个处理效应：
\[
\mathrm{Installs}_{it}
 =\sum_{i=0}^{N}\sum_{t=0}^{T}\tau_{it}D_{it}+\gamma_i+\theta_t+e_{it}.
\]
参数太多。改为给每个组 \(g\) 在每个时点分配一个效应：
\[
\mathrm{Installs}_{it}
 =\sum_{g=0}^{G}\sum_{t=0}^{T}\tau_{gt}D_{it}+\gamma_i+\theta_t+e_{it}.
\]
这里被标为更好的版本。进一步删去处理前系数和对照组系数，以减少参数：
\[
\mathrm{Installs}_{it}
 =\sum_{g=q}^{G}\sum_{t\geq g}^{T}\tau_{gt}D_{it}+\gamma_i+\theta_t+e_{it}.
\]
将 \(g<q\) 的队列定义为对照。
\notearrow{blue}{需要手动创建队列dummy，例如 \((0,1),(1,0),\ldots\)}{队列与时间交互的处理项。}\begin{reviewnote}{original}{B2-009}{需要显式写出队列与时间指标}{Chapter24-25.pdf，第4页}
对观测 \((i,s)\)，所要表达的展开应为
\[
Y_{is}=\alpha_i+\theta_s+
\sum_{g\in\mathcal G}\sum_{t\geq g}\tau_{g,t}
\mathbf1\{G_i=g\}\mathbf1\{s=t\}+e_{is}.
\]
队列指标与日历时间指标不可省略。这里修正的是公式记号，不是识别假设。
\par\smallskip{\small\textit{参考：} \href{https://matheusfacure.github.io/python-causality-handbook/24-The-Diff-in-Diff-Saga.html}{原始资料1}.}
\end{reviewnote}
\begin{reviewnote}{qualification}{B2-010}{对照队列必须提供未处理比较}{Chapter24-25.pdf, pp. 4}
若 \(g\) 是实际的采用日期，较早队列可能已经受处理，不能自动当作未处理对照。应采用从未处理单位，或在相关比较窗口内尚未处理的单位，并配合相应的平行趋势与无预期效应假设。教材示例把观测窗口结束后才采用的队列标为对照组。
\par\smallskip{\small\textit{参考：} \href{https://matheusfacure.github.io/python-causality-handbook/24-The-Diff-in-Diff-Saga.html}{原始资料1}.}
\end{reviewnote}

仍然存在的限制：
\begin{enumerate}
\item 效果会随时间变化，之前没有考虑这一点。
\item DiD和TWFE要满足平行趋势假设；这个限制很强。
\end{enumerate}
\notearrow{black}{DiD与TWFE}{需要相应的平行趋势限制。}

\section{第25章：合成双重差分}
\subsection{面板数据的表示}
\sourcepages{Chapter24-25.pdf, p. 5.}
多个单位、多个时期，受处理与未处理观测混合在一起。在 \(D\) 中，\textbf{行为时间，列为单位}：
\[
D=\begin{pmatrix}
0&0&0&\cdots&0&0\\
0&0&0&\cdots&0&0\\
\vdots&\vdots&\vdots&&\vdots&\vdots\\
0&0&0&\cdots&1&1\\
0&0&0&\cdots&1&1
\end{pmatrix}.
\]
最后几列对应受处理单位。
\notearrow{black}{采用后出现1的列}{受处理单位} 笔记称其为\emph{一次性处理}，并用星号强调：单位一旦受处理，就不再回到未处理状态。
\begin{reviewnote}{qualification}{B2-011}{四块矩阵要求共同的处理区块}{Chapter24-25.pdf, pp. 5}
图中的四块表示要求所纳入的受处理单位共享一个处理前／处理后的划分。若采用时间交错，应像后面的笔记那样，先构造有效的队列或时期子面板。
\par\smallskip{\small\textit{参考：} \href{https://arxiv.org/pdf/1812.09970}{原始资料1}, \href{https://matheusfacure.github.io/python-causality-handbook/25-Synthetic-Diff-in-Diff.html}{原始资料2}.}
\end{reviewnote}


\paragraph{表示1：合成控制。}
将上面的矩阵概括为四块：
\[
D=\begin{pmatrix}0&0\\0&1\end{pmatrix},\qquad
Y=\begin{pmatrix}
Y_{\mathrm{pre,co}}&Y_{\mathrm{pre,tr}}\\
Y_{\mathrm{post,co}}&Y_{\mathrm{post,tr}}
\end{pmatrix}.
\]
pre表示处理前，post表示处理后，co表示对照组，tr表示处理组。
\paragraph{表示2：DiD。}
用长表表示：各列记录单位、时间、结果、是否属于处理组，以及是否已经开始接受处理。
可观测结果为
\[
Y=\begin{pmatrix}
Y(0)_{\mathrm{pre,co}}&Y(0)_{\mathrm{pre,tr}}\\
Y(0)_{\mathrm{post,co}}&Y(1)_{\mathrm{post,tr}}
\end{pmatrix},
\]
并且
\[
\operatorname{ATT}=Y(1)_{\mathrm{post,tr}}
-\underbrace{Y(0)_{\mathrm{post,tr}}}_{\text{缺失}}.
\]

\subsection{回顾DiD与合成控制}
\sourcepages{Chapter24-25.pdf，第6、7页；第7页补充两幅比较示意图。}
DiD回归为
\[
Y_{it}=\beta_0+\beta_1\mathrm{Post}_t+\beta_2\mathrm{Treated}_i
+\beta_3\mathrm{Treated}_i\mathrm{Post}_t+e_{it}.
\]
其假设是
\[
E[Y(0)_{\mathrm{post,co}}-Y(0)_{\mathrm{pre,co}}]
=E[Y(0)_{\mathrm{post,tr}}-Y(0)_{\mathrm{pre,tr}}].
\]
即若没有处理，对照组和处理组具有相同趋势。

将估计量写成可与SDID比较的形式：
\[
\widehat\tau^{\mathrm{did}}
=\operatorname*{argmin}_{\mu,\alpha,\beta,\tau}
\left\{\sum_{i=1}^{N}\sum_{t=1}^{T}
[Y_{it}-(\mu+\alpha_i+\beta_t+\tau D_{it})]^2\right\}.
\]
笔记将它与去均值联系起来：
\begin{align*}
\widetilde Y_{it}
&=Y_{it}-\overline Y_i-\overline Y_t
=Y_{it}-T^{-1}\sum_{t=0}^{T}Y_{it}-N^{-1}\sum_{i=0}^{N}Y_{it},\\
\widetilde D_{it}
&=D_{it}-\overline D_i-\overline D_t
=D_{it}-T^{-1}\sum_{t=0}^{T}D_{it}-N^{-1}\sum_{i=0}^{N}D_{it}.
\end{align*}
这与合成控制很像，但没有权重。
\notearrow{black}{对 \(Y\) 和 \(D\) 去均值}{固定效应的最小二乘目标函数。}
\begin{reviewnote}{original}{B2-004}{再次出现的双重去均值遗漏}{Chapter24-25.pdf, p. 6}
应加回总体均值，并统一求和上下限；见前面的B2-004校正。
\par\smallskip{\small\textit{参考：} \href{https://matheusfacure.github.io/python-causality-handbook/24-The-Diff-in-Diff-Saga.html}{原始资料1}, \href{https://matheusfacure.github.io/python-causality-handbook/25-Synthetic-Diff-in-Diff.html}{原始资料2}.}
\end{reviewnote}

合成控制使用处理前结果估计权重：
\[
\widehat w^{\mathrm{sc}}
=\operatorname*{argmin}_{w}
\|\overline y_{\textcolor{blue}{\mathrm{pre,tr}}}-Y_{\mathrm{pre,co}}w_{\mathrm{co}}\|_2^2,
\qquad \sum_iw_i=1,\quad w_i>0\quad\forall i.
\]
\begin{reviewnote}{original}{B2-012}{应允许零对照权重}{Chapter24-25.pdf, pp. 6, 8, 10}
标准单纯形约束为 \(w_i\geq0\)、\(\sum_{i\in\mathrm{co}}w_i=1\)；同理，\(\lambda_t\geq0\)、\(\sum_{t\in\mathrm{pre}}\lambda_t=1\)。允许零权重。开单纯形 \(w_i>0\) 可能排除最优解。第8、10页的手写约束还遗漏了求和号：并非每个权重都等于1。下文公式恢复了权重总和为1的本意。
\par\smallskip{\small\textit{参考：} \href{https://matheusfacure.github.io/python-causality-handbook/25-Synthetic-Diff-in-Diff.html}{原始资料1}, \href{https://arxiv.org/pdf/1812.09970}{原始资料2}.}
\end{reviewnote}

\(Y_{\mathrm{pre,co}}\) 是 \(T_{\mathrm{pre}}\times N_{\mathrm{co}}\) 的结果矩阵：列为单位，行为时期。向量 \(w_{\mathrm{co}}\) 的大小为 \(N_{\mathrm{co}}\times1\)。\(\overline y_{\mathrm{pre,tr}}\) 的每个分量，是某个处理前时点上处理组单位的平均结果。

\paragraph{我的批注。}
\inkblue{给处理组拟合一个合成对照，类似线性回归。}处理后，用对照组结果与已经拟合的权重计算：
\[
y_{\mathrm{sc}}=Y_{\mathrm{co}}\widehat w^{\mathrm{sc}}.
\]
这给出了 \(Y(0)_{\mathrm{post,tr}}\) 的良好估计，因此
\[
\widehat\tau=\overline y_{\mathrm{post,tr}}-\overline y_{\mathrm{post,sc}}.
\]
也可以改写为
\[
\widehat\tau^{\mathrm{sc}}
=\operatorname*{argmin}_{\beta,\tau}
\left\{\sum_{i=1}^{N}\sum_{t=1}^{T}
(Y_{it}-\beta_t-\tau D_{it})^2\widehat w_i^{\mathrm{sc}}\right\}.
\]
处理前的拟合给出对照组权重；每个处理组单位的权重为 \(1/N_{\mathrm{tr}}\)。\notearrow{blue}{没有总体截距 \(\mu\)，也没有单位固定效应 \(\alpha_i\)}{上面的SC加权回归目标函数。}

\paragraph{左页的图形直觉（第7页）。}
DiD图将轨迹平移，构造平行比较；手写批注说，这在图示例子中放大了效果。SC图则从处理前数据得到权重，形成一个合成的、未受处理的处理组。处理后继续用同一组权重计算结果，并以两个处理后结果之差估计效果。图中曲线在处理前贴合，处理后分开。
\notearrow{black}{处理前拟合}{权重 \(\widehat w\)}
\notearrow{black}{处理后的对照结果与同一组权重}{合成的处理后结果，再与处理组结果求差。}
\begin{center}
\begin{tikzpicture}[x=0.75cm,y=0.6cm,>=stealth]
\draw[->] (0,0)--(5,0) node[right]{时间};
\draw[->] (0,0)--(0,3);
\draw[dashed] (2.5,0)--(2.5,3);
\draw (0.4,2.8)--(4.4,1.1);
\draw (0.4,2.1)--(4.4,0.25);
\draw[dashed] (0.4,2.1)--(4.4,0.85);
\node at (2.2,-0.7){DiD：平移趋势};
\begin{scope}[xshift=7.1cm]
\draw[->] (0,0)--(5,0) node[right]{时间};
\draw[->] (0,0)--(0,3);
\draw[dashed] (2.5,0)--(2.5,3);
\draw (0.4,2.65) .. controls (1.7,2.7) and (2,2.1) .. (2.5,1.8)
.. controls (3,1.5) and (3.5,1.3) .. (4.4,1.3);
\draw (0.4,2.55) .. controls (1.7,2.6) and (2,2.2) .. (2.5,1.75)
.. controls (2.8,0.6) and (3.2,.45) .. (4.4,.45);
\draw[<->] (4.6,.45)--(4.6,1.3);
\node at (2.4,-0.7){SC：加权轨迹};
\end{scope}
\end{tikzpicture}

{\small 按原笔记草图重建的示意图。}
\end{center}


\subsection{将两者结合：合成DiD}
\sourcepages{Chapter24-25.pdf，第8页与第10页上、中部。}
\inkred{\(\star\)} 星号标题将合成控制的加权目标与DiD固定效应目标结合起来：
\[
\widehat\tau^{\mathrm{sdid}}
=\operatorname*{argmin}_{\mu,\alpha,\beta,\tau}
\left\{\sum_{i=1}^{N}\sum_{t=1}^{T}
[Y_{it}-(\mu+\alpha_i+\beta_t+\tau D_{it})]^2
\underbrace{\widehat w_i^{\mathrm{sdid}}}_{\textcolor{red}{\text{单位权重}}}
\underbrace{\widehat\lambda_t^{\mathrm{sdid}}}_{\textcolor{red}{\text{时间权重}}}\right\}.
\]
\paragraph{时间权重。}
用对照组处理前结果拟合其处理后平均结果：
\[
\widehat\lambda^{\mathrm{sdid}}
=\operatorname*{argmin}_{\lambda}
\|\overline y_{\textcolor{blue}{\mathrm{post,co}}}
-(\lambda_{\mathrm{pre}}Y_{\mathrm{pre,co}}+\lambda_0)\|_2^2,
\qquad \sum_t\lambda_t=1,\quad\lambda_t>0.
\]
\begin{itemize}
\item \(Y_{\mathrm{pre,co}}\) 为 \(T_{\mathrm{pre}}\times N_{\mathrm{co}}\) 矩阵，行为时间、列为单位。
\item \(\overline y_{\mathrm{post,co}}\) 是 \(1\times N_{\mathrm{co}}\) 行向量，每个分量是某个对照单位在处理后各时期的平均结果。
\item \(\lambda_{\mathrm{pre}}\) 是 \(1\times T_{\mathrm{pre}}\) 行向量，为每个处理前时期指定一个权重；\(\lambda\) 是时间权重。
\end{itemize}
第8页底部用于替代划掉文字的矩阵草图，沿时间维度相乘：在矩阵上方添加一行1表示截距，再用加权的处理前各行拟合 \(\overline y_{\mathrm{post,co}}\)，从而得到 \(\lambda_t\)。
\begin{center}
\begin{tikzpicture}[>=stealth,node distance=22mm]
\node (x) {\(\begin{pmatrix}\mathbf1^{\mathsf T}\\Y_{\mathrm{pre,co}}\end{pmatrix}\)};
\node[right=of x] (fit) {\((\lambda_0,\lambda_{\mathrm{pre}})\)};
\node[right=of fit,blue] (y) {\(\overline y_{\mathrm{post,co}}\)};
\draw[->] (x)--node[above,font=\small]{时间拟合}(fit);
\draw[->] (fit)--(y);
\end{tikzpicture}
\end{center}

\paragraph{单位权重。}
第10页顶部的小图用 \(\lambda\) 连接“对照／处理前”与“对照／处理后”，用 \(w\) 连接“对照／处理前”与“处理组／处理前”，从而得到相应的处理组／处理后比较。
\begin{center}
\begin{tikzpicture}[>=stealth,node distance=14mm and 28mm,blue]
\node (a) {对照／处理前}; \node[right=of a] (b) {对照／处理后};
\node[below=of a] (c) {处理组／处理前}; \node[below=of b] (d) {处理组／处理后};
\draw[->] (a)--node[above]{\(\lambda\)}(b); \draw[->] (a)--node[left]{\(w\)}(c);
\draw[->] (b)--(d); \draw[->] (c)--(d);
\end{tikzpicture}
\end{center}
\begin{align*}
\widehat w^{\mathrm{sdid}}
=\operatorname*{argmin}_{w}\ \big\{
&\|\overline y_{\textcolor{blue}{\mathrm{pre,tr}}}
-(Y_{\mathrm{pre,co}}w_{\mathrm{co}}+w_0)\|_2^2\\
&+\zeta^2T_{\mathrm{pre}}\|w_{\mathrm{co}}\|_2^2\big\},
\qquad \sum_iw_i=1,\quad w_i>0,\\
\zeta&=(N_{\mathrm{tr}}T_{\mathrm{post}})^{1/4}
\widehat\sigma(\Delta_{it}),\\
\Delta_{it}&=Y_{it}-Y_{i,t-1}.
\end{align*}
其中 \(\Delta_{it}\) 是一阶差分，\(\widehat\sigma(\Delta_{it})\) 是其标准差。

\paragraph{我的批注。}
\notearrow{red}{\(L_2\)惩罚让非零权重分布得更广}{\(\zeta^2T_{\mathrm{pre}}\|w_{\mathrm{co}}\|_2^2\)}
\notearrow{blue}{与之前SC相比：（1）有截距；（2）有 \(L_2\) 惩罚}{SDID的单位权重拟合。} \notearrow{blue}{有截距}{\(w_0\)}
原稿的实现说明写的是：在\emph{\(\overline y_{\mathrm{pre,tr}}\) 前添加一行1}作为截距，再做回归。
\begin{reviewnote}{original}{B2-013}{截距应放在设计矩阵中}{Chapter24-25.pdf, pp. 10}
在单位权重拟合中，\(\overline y_{\mathrm{pre,tr}}\) 是长度为 \(T_{\mathrm{pre}}\) 的向量，设计矩阵应为
\[
X=[\,\mathbf1_{T_{\mathrm{pre}}}\ \ Y_{\mathrm{pre,co}}\,].
\]
应在预测变量中添加一\emph{列}1，而不是给结果向量加一行。拟合 \((w_0,w_{\mathrm{co}})\) 时，只约束并惩罚对照单位的权重，不惩罚截距。相比之下，第8页时间权重采用行向量表示，因此在 \(Y_{\mathrm{pre,co}}\) 上方添加一行1是正确的。
\par\smallskip{\small\textit{参考：} \href{https://matheusfacure.github.io/python-causality-handbook/25-Synthetic-Diff-in-Diff.html}{原始资料1}, \href{https://arxiv.org/pdf/1812.09970}{原始资料2}.}
\end{reviewnote}
\begin{reviewnote}{qualification}{B2-014}{应使用对照组处理前数据估计尺度}{Chapter24-25.pdf, pp. 10}
标准SDID使用对照单位在处理前时期的结果一阶差分来估计 \(\widehat\sigma\)。若使用处理组的处理后结果，处理效应可能影响正则化尺度。笔记中的四次方根因子 \((N_{\mathrm{tr}}T_{\mathrm{post}})^{1/4}\)（指数为正）是正确的。
\par\smallskip{\small\textit{参考：} \href{https://arxiv.org/pdf/1812.09970}{原始资料1}.}
\end{reviewnote}


处理组单位与处理后时期分别使用均匀权重 \(1/N_{\mathrm{tr}}\) 和 \(1/T_{\mathrm{post}}\)。\textbf{我的直觉：}SDID相比SC，更关注干预近期起效果的数据。\begin{reviewnote}{qualification}{B2-015}{SDID直觉中的字迹仍有歧义}{Chapter24-25.pdf，第10页}
上面按时间接近程度理解是可能的，但原句中间几个字仍不清楚。数学上，时间权重使加权后的对照组处理前结果能够预测其处理后结果；它不会依据观测与某个未知处理效应有多接近来选数据，也不一定只选最近的时期。
\par\smallskip{\small\textit{参考：} \href{https://matheusfacure.github.io/python-causality-handbook/25-Synthetic-Diff-in-Diff.html}{原始资料1}, \href{https://arxiv.org/pdf/1812.09970}{原始资料2}.}
\end{reviewnote}


\subsection{时间效应异质性与交错采用}
\sourcepages{Chapter24-25.pdf，第10页下部，随后接第9页补充。}
SDID可以方便地估计各个时期的效应：重复运行即可。下面手写矩阵中0和1的下标表示时期。
\paragraph{\inkred{时间效应异质性。}}
例子为
\[
D=\begin{pmatrix}
0_1&0_1&0_1\\0_2&0_2&0_2\\0_3&0_3&1_3\\0_4&0_4&1_4
\end{pmatrix}.
\]
对整个矩阵使用SDID，可得第3、4期的平均ATT。
\notearrow{black}{完整矩阵 \(D\)}{SDID：第3、4期的平均ATT。} 若要分别估计这两期，则用
\[
D_1=\begin{pmatrix}0_1&0_1&0_1\\0_2&0_2&0_2\\0_3&0_3&1_3\end{pmatrix},
\qquad
D_2=\begin{pmatrix}0_1&0_1&0_1\\0_2&0_2&0_2\\0_4&0_4&1_4\end{pmatrix}.
\]

\paragraph{\inkred{交错采用：左页补充。}}
笔记给出
\[
D=\begin{pmatrix}
0_1&0_1&0_1\\0_2&0_2&0_2\\0_3&0_3&1_3\\0_4&1_4&1_4
\end{pmatrix}.
\]
第一种拆分为
\[
D_1=\begin{pmatrix}0_1&0_1&0_1\\0_2&0_2&0_2\\0_4&1_4&1_4\end{pmatrix},
\qquad
D_2=\begin{pmatrix}0_1&0_1&0_1\\0_2&0_2&0_2\\0_3&0_3&1_3\end{pmatrix}.
\]
将 \(D_1\) 的ATT赋予权重 \(2/3\)，将 \(D_2\) 的ATT赋予权重 \(1/3\)。
\notearrow{black}{按时期拆分原处理矩阵}{第4期有两个受处理单元，第3期有一个：权重为 \(2/3\)、\(1/3\)。} 也可以按受处理队列拆分：
\[
D_1=\begin{pmatrix}0_1&0_1\\0_2&0_2\\0_3&0_3\\0_4&1_4\end{pmatrix},
\qquad
D_2=\begin{pmatrix}0_1&0_1\\0_2&0_2\\0_3&1_3\\0_4&1_4\end{pmatrix}.
\]
下面的例子写道：加州受处理12年，其他单位8年 \(\times3\)。加权计算为
\[
\mathrm{ATT}_0\frac{12}{36}
+\left(\sum\mathrm{ATT}\right)\frac13\frac{3\times8}{36}.
\]
这里保留了作者按处理持续时间加权的例子。
\begin{reviewnote}{qualification}{B2-017}{受处理单元的权重决定聚合目标}{Chapter24-25.pdf, pp. 9, 10}
权重 \(2/3,1/3\) 按受处理的单位—时期单元计数。同理，\(12/36\) 与 \(24/36\) 按已观测的处理持续时间，平均加州与其他三个单位的效应，每个单位、每个时期的权重相等。其他目标总体可能需要不同聚合权重。各子面板都应使用真正未受处理的对照观测。
\par\smallskip{\small\textit{参考：} \href{https://matheusfacure.github.io/python-causality-handbook/25-Synthetic-Diff-in-Diff.html}{原始资料1}, \href{https://arxiv.org/pdf/1812.09970}{原始资料2}.}
\end{reviewnote}


\subsection{安慰剂方差估计}
\sourcepages{Chapter24-25.pdf, p. 11.}
笔记描述SDID相比SC有更小的误差棒。SDID同时利用结果在时间和单位之间的变化，使置信区间改善，而且适用于只有一个受处理单位的情形。
\begin{reviewnote}{qualification}{B2-018}{精度并非总能提高}{Chapter24-25.pdf, pp. 11}
SDID在适当设计下可以提高精度和稳健性，但其有效标准误并不保证总比SC小。估计权重本身也带来不确定性，因此普通加权回归输出的标准误并不自动等于SDID的标准误。
\par\smallskip{\small\textit{参考：} \href{https://matheusfacure.github.io/python-causality-handbook/25-Synthetic-Diff-in-Diff.html}{原始资料1}, \href{https://arxiv.org/pdf/1812.09970}{原始资料2}.}
\end{reviewnote}


轮流假设每个对照单位接受了处理，用SDID估计其安慰剂ATT。笔记写为
\[
\widehat V_{\tau}^{\mathrm{placebo}}
=B^{-1}\sum_{b=1}^{B}
\left(\widehat\tau^{(b)}-\overline\tau^{(b)}\right)^2,
\qquad
\tau\in\widehat\tau^{\mathrm{sdid}}
\pm z_{\alpha/2}\sqrt{\widehat V_\tau}.
\]
\notearrow{black}{安慰剂ATT估计值的方差}{SDID安慰剂方差 \(\widehat V_\tau\)，再构造置信区间。} 括号中批注：和SC一样。
\begin{reviewnote}{qualification}{B2-019}{定义安慰剂均值与区间约定}{Chapter24-25.pdf, pp. 11}
对 \(B\) 个安慰剂估计，定义 \(\overline\tau^{\mathrm{pl}}=B^{-1}\sum_b\widehat\tau^{(b)}\)，并令
\[
\widehat V^{\mathrm{pl}}=B^{-1}\sum_b(\widehat\tau^{(b)}-\overline\tau^{\mathrm{pl}})^2.
\]
若定义 \(z_p=\Phi^{-1}(p)\)，正态近似区间应为 \(\widehat\tau\pm z_{1-\alpha/2}\sqrt{\widehat V^{\mathrm{pl}}}\)。手写的 \(z_{\alpha/2}\) 也可能约定为正的上尾临界值，但原稿未说明。置信区间是具有覆盖率声明的随机集合，并不保证每次实现出的区间都包含 \(\tau\)。
\par\smallskip{\small\textit{参考：} \href{https://arxiv.org/pdf/1812.09970}{原始资料1}.}
\end{reviewnote}
\begin{reviewnote}{qualification}{B2-020}{单个处理单位的安慰剂推断需要假设}{Chapter24-25.pdf, pp. 11}
安慰剂方法从对照组借用噪声信息。它要求不同单位的噪声可比较（SDID论文的这一论证要求跨单位同方差），有适当的对照样本，并能合理使用相应的分布近似。若处理单位的噪声变异不同，方法可能失效。多个单位受处理时，每次安慰剂分配应保持相同的处理组规模。
\par\smallskip{\small\textit{参考：} \href{https://arxiv.org/pdf/1812.09970}{原始资料1}.}
\end{reviewnote}


\section{附录：去偏与预测}
\sourcepages{Chapter24-25.pdf, p. 12.}
页首把非随机数据的处理与\textbf{倾向得分}、\textbf{正交化}联系起来。
\subsection{通过正交化去偏}
\notearrow{blue}{与第22章相同}{正交化／对变量取残差。} 所有特征包含在 \(X\) 中。
\begin{enumerate}
\item 取第一组特征，拟合
\[
y_i=\theta_0+\theta_1X_{1i}+e_i.
\]
得到残差
\[
y_i^*=y_i-(\widehat\theta_0+\widehat\theta_1X_{1i}).
\]
\item 用第一组特征预测第二组：
\[
X_{2i}=\gamma_0+\gamma_1X_{1i}+e_i.
\]
得到残差
\[
X_{2i}^*=X_{2i}-(\widehat\gamma_0+\widehat\gamma_1X_{1i}).
\]
\item 用特征残差回归结果残差：
\[
y_i^*=\beta_0+\beta_2X_{2i}^*+e_i.
\]
\end{enumerate}

\subsection{通过倾向得分去偏}
\sourcepages{Chapter24-25.pdf, p. 13.}
\notearrow{blue}{见第11章}{倾向得分加权。} 对每个单位加权。二元处理时，
\[
w_i=\frac{T}{P(T=1\mid X)}+\frac{1-T}{1-P(T=1\mid X)}.
\]
原稿把下式标为\emph{连续}：
\[
w_i=\sum_{k=0}^{K}\frac{\mathbf1\{T=k\}}{P(T=k\mid X)}.
\]
\begin{reviewnote}{original}{B2-021}{所写扩展适用于离散类别}{Chapter24-25.pdf, pp. 13}
这个表达式用于取值为 \(0,\ldots,K\) 的离散处理。对于真正连续的剂量，常见连续模型在某个单点上满足 \(P(T=k\mid X)=0\)；广义倾向得分方法使用条件密度，并需要另行明确剂量—反应目标量与估计方法。
\par\smallskip{\small\textit{参考：} \href{https://matheusfacure.github.io/python-causality-handbook/Debiasing-with-Propensity-Score.html}{原始资料1}.}
\end{reviewnote}


\paragraph{我的直觉。}
\begin{inkblock}{blue}
受处理概率很高却未受处理的单位，本质上与处理组很像，但观测上属于对照组，可用于估计 \(Y_0\mid T=1\)。
\end{inkblock}\begin{reviewnote}{qualification}{B2-022}{ATE与ATT的权重不同}{Chapter24-25.pdf, pp. 13}
在一致性、条件可交换性与正值性成立时，上面的逆概率权重对应ATE比较。若目标为ATT，处理组权重应为1，对照组权重为 \(e(X)/(1-e(X))\)，并作适当归一化。高倾向得分的对照单位有助于代表处理组的协变量分布，但它并不真的就是受处理单位。
\par\smallskip{\small\textit{参考：} \href{https://matheusfacure.github.io/python-causality-handbook/Debiasing-with-Propensity-Score.html}{原始资料1}.}
\end{reviewnote}

没有倾向得分时，可以用机器学习估计接受处理的概率。一个简短数值例子提到20\%、0.2，以及“收到邮件”记为 \(1\)；这一行的其余字迹不清。
\begin{reviewnote}{qualification}{B2-023}{邮件例子仍有字迹无法辨认}{Chapter24-25.pdf, pp. 13}
保留可辨认的数值，不从缺失文字中推测额外的概率计算或因果结论。
\end{reviewnote}


\paragraph{倾向得分的弱点。}
\begin{enumerate}
\item 倾向得分很小时（\(e<0.05\)），权重过大，相当于该观测重复出现太多次，对去偏后的数据集影响很大。笔记提到人工裁剪。\inkblue{\textbf{高方差。}}
\item \textbf{正值性或共同支持：}处理单位与控制单位的特征之间必须有足够重叠。重叠过少时，用倾向得分消除相关性的效果不好。
\end{enumerate}
\begin{inkblock}{red}
人工干预不太好，可能引入过多主观偏差；不要外推到正值性不成立的人口子集。
\end{inkblock}\begin{reviewnote}{qualification}{B2-024}{大权重与样本裁剪}{Chapter24-25.pdf, pp. 13}
二元ATE加权中，\(e\) 很小时膨胀的是处理单位的 \(1/e\) 权重；\(e\) 接近1时膨胀的是对照单位的 \(1/(1-e)\) 权重。针对重叠不足，可以按事先规定的原则裁剪样本，但这会改变目标总体，必须报告。良好的倾向得分预测本身并不能证明因果可交换性。
\par\smallskip{\small\textit{参考：} \href{https://matheusfacure.github.io/python-causality-handbook/Debiasing-with-Propensity-Score.html}{原始资料1}.}
\end{reviewnote}


\subsection{当预测失效时}
\sourcepages{Chapter24-25.pdf, p. 14.}
预测与因果。这一页只有标题，下面没有更多笔记。
